\chapter*{Введение}                         % Заголовок
\addcontentsline{toc}{chapter}{Введение}    % Добавляем его в оглавление

\subsection*{Актуальность темы исследования}
Актуальность исследования промышленного цикла в капиталистической экономике определяется объективной ролью циклических колебаний в воспроизводстве общественного капитала, поскольку смена кризиса, депрессии, оживления и подъема представляет собой не случайное нарушение нормального хода хозяйственной жизни, а закономерную форму движения капиталистического производства, посредством которой обостряются, проявляются и временно разрешаются противоречия между производством и потреблением, производством и обращением, накоплением и реализацией, общественным характером труда и частнокапиталистической формой присвоения его результатов. В промышленном цикле обнаруживается внутренняя необходимость периодического восстановления нарушенных пропорций общественного воспроизводства через обесценение капитала, сокращение производства, разрушение части производительных сил, перераспределение собственности и создание новых условий для последующего возобновления накопления.

Промышленный цикл представляет собой верхушку айсберга всего капиталистического способа производства в его исторической ретроспективе, поскольку в последовательности его фаз в концентрированном виде отражаются глубинные закономерности возникновения, становления и развития капиталистических производственных отношений, эволюция форм собственности, кредита, конкуренции, монополий, мирового рынка и государственного вмешательства в хозяйственную жизнь. Исследование цикла тем самым позволяет перейти от внешне наблюдаемых колебаний объемов производства, инвестиций, занятости и цен к анализу лежащих в их основе отношений капитала и наемного труда, производства прибавочной стоимости, формирования нормы прибыли, перелива капитала между отраслями и исторически изменяющихся форм подчинения общественного производства потребностям самовозрастания стоимости.

На протяжении почти двух столетий капиталистическая экономика с определенной регулярностью воспроизводит условия кризисов общего перепроизводства, которые по мере углубления общественного разделения труда, концентрации и централизации капитала, интернационализации производства, расширения кредита и формирования глобальных производственно-финансовых связей приобретают все более широкий масштаб и оказывают все более разрушительное воздействие на производство, занятость, доходы трудящихся и устойчивость общественного воспроизводства. При этом кризисы возникают не вследствие абсолютного избытка произведенных благ по отношению к общественным потребностям, а вследствие их относительного перепроизводства по отношению к ограниченному платежеспособному спросу, обусловленному антагонистическим характером распределения, стремлением капитала к безграничному расширению производства и одновременным сдерживанием потребления основной массы населения.

Циклическая динамика пронизывает практически все сферы хозяйственной жизни, охватывая не только промышленное производство, но и инвестиционные процессы, строительство, транспорт, торговлю, кредит, денежное обращение, финансовые рынки, занятость, заработную плату, государственные финансы и внешнеэкономические связи, вследствие чего промышленный цикл выступает всеобщей формой движения капиталистического воспроизводства, а не изолированным явлением отдельных отраслей. Смена фаз цикла сопровождается изменением спроса на рабочую силу, степени загрузки производственных мощностей, массы и нормы прибыли, объемов основного и оборотного капитала, интенсивности кредитования, структуры совокупного общественного продукта и пропорций между подразделениями общественного производства, что придает исследованию цикличности фундаментальное значение для понимания всей системы хозяйственных взаимосвязей капиталистической экономики.

Особую актуальность теме придает ограниченность распространенных немарксистских концепций экономического цикла, которые, абстрагируясь от капиталистических производственных отношений, нередко объясняют кризисы внешними шоками, ошибками денежно-кредитной политики, изменениями инвестиционных ожиданий, несовершенством информации, технологическими импульсами либо иррациональностью поведения хозяйствующих субъектов и тем самым мистифицируют действительный процесс воспроизводства, подменяя исследование его внутренних противоречий описанием отдельных поверхностных форм их проявления. В этих условиях возникает необходимость углубленного комплексного исследования промышленного цикла на основе марксистского учения, позволяющего рассматривать кризис как имманентный момент движения капитала, связать производство прибавочной стоимости с условиями ее реализации и капитализации и раскрыть единство процессов производства, обращения, распределения и потребления в рамках воспроизводства совокупного общественного капитала.

Необходимость исследования промышленного цикла обусловлена также исторической изменчивостью капитализма, поскольку развитие монополий и финансового капитала, усиление роли государства, формирование транснациональных корпораций, глобализация производственных цепочек, финансиализация, цифровизация и ускорение технологического обновления существенно изменяют конкретные механизмы, продолжительность, глубину и формы проявления циклических колебаний, не устраняя при этом их объективной основы. Современная циклическая динамика характеризуется тесным переплетением промышленных, финансовых, долговых, структурных, технологических и внешнеэкономических противоречий, а также возможностью переноса кризисных импульсов между странами и отраслями через мировую торговлю, движение капитала, кредитные отношения и глобальные цепочки стоимости, что требует рассматривать промышленный цикл как исторически конкретную форму движения мирового капиталистического хозяйства.

Таким образом, актуальность темы исследования обусловлена, с одной стороны, фундаментальным значением промышленного цикла для понимания закономерностей капиталистического воспроизводства, а с другой – необходимостью теоретического осмысления новых форм проявления экономической цикличности в условиях современного капитализма, когда государственное регулирование, развитие кредита, финансовых рынков и транснациональных производственных структур модифицируют внешние формы кризисов, но не устраняют лежащих в их основе противоречий накопления капитала. Исследование промышленного цикла на основе диалектического и историко-материалистического метода позволяет раскрыть единство производства, обращения, распределения и потребления, установить внутреннюю связь между накоплением капитала, изменением нормы прибыли, обновлением основного капитала и кризисными нарушениями воспроизводства, определить объективные и исторические пределы государственного антициклического регулирования, а также создать теоретико-методологическую основу для дальнейшего развития марксистской политической экономии и конкретного анализа динамики капиталистического хозяйства.

\subsection*{Степень разработанности научной проблемы}
В развитии и становлении теоретических концепций экономического цикла можно выделить три основные линии, различающиеся по своим методологическим основаниям, пониманию природы капиталистического хозяйства и трактовке источников циклического движения: марксистскую политическую экономию, буржуазную экономическую теорию и совокупность эклектических концепций, занимающих промежуточное положение между ними. Такое разграничение основывается не столько на формальной принадлежности отдельных авторов к той или иной научной школе, сколько на внутренней логике их исследований, используемом понятийном аппарате и решении вопроса о том, выводится ли промышленный цикл из имманентных противоречий капиталистического способа производства либо объясняется воздействием отдельных внешних или функционально обособленных факторов.

Фундаментальные основы марксистской теории промышленного цикла были заложены К. Марксом, Ф. Энгельсом, В. И. Лениным и И. В. Сталиным, в трудах которых кризисы общего перепроизводства были раскрыты как закономерная форма движения и временного разрешения противоречий капиталистического воспроизводства, связанных с производством и реализацией прибавочной стоимости, накоплением и обесценением капитала, оборотом основного капитала, развитием кредита, тенденцией нормы прибыли к понижению, концентрацией и централизацией производства, формированием финансового капитала и мирового хозяйства. Дальнейшее развитие данных положений было осуществлено советскими экономистами П. С. Максаковским, Л. А. Мендельсоном, В. Г. Фигурновым, Е. С. Варгой, И. А. Трахтенбергом и другими исследователями, которые провели более глубокий системный анализ структуры и фаз промышленного цикла, механизма кризиса, роли обновления основного капитала, кредита, денежного обращения и межотраслевых диспропорций, а также разработали историческую периодизацию кризисов и исследовали изменение циклического движения в условиях монополистического и государственно-монополистического капитализма; в современных работах М. Робертса, Г. Карчеди, В. Калюжного и ряда других авторов марксистские категории применяются к анализу новейшего статистического материала, предлагаются различные методики расчета органического строения капитала, нормы прибыли, степени эксплуатации труда и международного перераспределения прибавочной стоимости, а результаты эмпирических исследований подтверждают основные положения марксистской политической экономии о росте органического строения капитала и действии закона тенденции нормы прибыли к понижению.

Буржуазные теории экономического цикла исторически начали формироваться ранее целостной марксистской концепции кризисов и до настоящего времени сохраняют господствующее положение в экономическом мейнстриме, системе экономического образования и практике государственного регулирования, однако представители различных и относительно самостоятельных школ существенно расходятся в определении природы и причин циклических колебаний. В большинстве случаев экономический цикл сводится ими к действию одного преобладающего фактора или ограниченной совокупности факторов – изменениям денежного предложения и процентной ставки, кредитной экспансии, инвестиционным колебаниям, недостаточности совокупного спроса, технологическим шокам, инновациям, несовершенству информации, негибкости цен либо изменению ожиданий хозяйствующих субъектов, – вследствие чего отдельные действительные моменты капиталистической динамики отрываются от целостного воспроизводства общественного капитала; значительный вклад в конкретизацию буржуазной теории цикла внес Й. Шумпетер, связавший циклические колебания с неравномерным внедрением инноваций и способствовавший распространению представления о сосуществовании циклов различной продолжительности, среди которых наибольшую известность получили краткосрочные циклы К. Китчина, среднесрочные циклы К. Жугляра, строительные и инфраструктурные ритмы С. Кузнеца, а также теория реального делового цикла Ф. Кидланда и Э. Прескотта.

Промежуточное положение между марксистской и буржуазной политической экономией занимают эклектические теории, которые в различной степени соединяют методологически несовместимые элементы двух названных направлений либо внешне воспроизводят марксистский понятийный аппарат, фактически заимствуя содержательные положения теорий недопотребления, автономного инвестиционного спроса, денежно-кредитной нестабильности, внешних рынков или технологического детерминизма. К данной группе в рамках принятой классификации могут быть отнесены теория длинных волн Н. Д. Кондратьева, отдельные положения работ Р. Люксембург, Н. И. Бухарина, Р. Гильфердинга и некоторых представителей так называемого легального марксизма Российской империи, поскольку в их концепциях марксистская терминология нередко сочетается с положениями, затрудняющими последовательное выведение кризисов из внутренних противоречий производства и накопления капитала; подобные построения приводят к перенесению причин кризиса из сферы капиталистического производства в сферу обращения, кредита, потребления или внешнего рынка, вследствие чего они неоднократно становились предметом полемики со стороны представителей последовательной марксистской политической экономии, выявлявших их теоретическую непоследовательность, оппортунистические тенденции и ограниченность при объяснении закономерного воспроизводства капиталистических кризисов.

Вопросы теоретической систематизации различных концепций экономического цикла были рассмотрены в работах К. Ю. Курилова и А. А. Куриловой, Д. В. Манушина, С. П. Аукуционека, Д. Бесоми, Х. Хагеманна, М. Бояновски и других исследователей, которыми предпринимались попытки классифицировать существующие теории в зависимости от предполагаемого источника колебаний, продолжительности цикла, характера передаточного механизма, отношения к экономическому равновесию, роли денежных и реальных факторов, а также эндогенного или экзогенного происхождения кризисных импульсов. Данные исследования имеют существенное значение для реконструкции истории экономической мысли и выявления преемственности между отдельными научными школами, однако предлагаемые в них классификации нередко строятся на одном доминирующем основании и потому не позволяют в полной мере отразить сложную внутреннюю структуру каждой теории, одновременно содержащей определенные представления о производстве, распределении, потреблении, кредите, техническом прогрессе, государственном регулировании и исторической направленности развития капитализма.

Вопросы математического моделирования циклических процессов в экономике получили разработку в трудах Ю. А. Кузнецова, А. А. Акаева, В. А. Петухова, Ю. Л. Плущевской, М. Калецкого, Р. Гудвина и других авторов, исследовавших возможности формального описания взаимодействия накопления, инвестиций, потребления, занятости, заработной платы, прибыли, технического прогресса и межотраслевых пропорций с помощью систем дифференциальных и разностных уравнений, моделей мультипликатора и акселератора, нелинейных динамических систем, имитационного моделирования и эконометрической обработки временных рядов. Особое значение имеют модели М. Калецкого, связывающие инвестиционные решения, прибыль и загрузку производственных мощностей, а также модель Р. Гудвина, описывающая циклическое взаимодействие доли заработной платы и уровня занятости, однако математическая строгость подобных построений сама по себе не обеспечивает адекватности их политико-экономического содержания, поскольку результаты моделирования непосредственно зависят от исходных теоретических предпосылок, выбора переменных и характера предполагаемых связей между ними.

Несмотря на значительный интерес научного сообщества к исследованию циклов в экономике, по-прежнему остается нерешенным ряд взаимосвязанных методологических проблем. Во-первых, существующие систематизации теорий цикла в большинстве случаев являются ограниченными и методологически неполноценными, поскольку их авторы стремятся разработать единый универсальный критерий разделения концепций – например, на денежные и реальные, эндогенные и экзогенные, равновесные и неравновесные, – тогда как более продуктивным представляется построение комплексной многоуровневой системы конкретных критериев, позволяющих сопоставлять позиции исследователей применительно к отдельным вопросам происхождения кризиса, механизма смены фаз, роли кредита, продолжительности цикла и исторических границ капитализма; во-вторых, основная масса современных работ исходит из преимущественно эмпирического или эконометрического анализа статистических показателей, уделяя недостаточное внимание конкретно-историческому, диалектико-материалистическому и историко-материалистическому исследованию циклических процессов, хотя статистические методы должны выступать средством проверки и конкретизации теоретических положений, а не заменять раскрытие внутренних противоречий капиталистического воспроизводства; в-третьих, большинство авторов чрезмерно концентрируется на отдельных сторонах циклической динамики – движении нормы прибыли, инвестициях, финансовой нестабильности, обновлении основного капитала или формах кризиса, – вследствие чего практически отсутствует комплексный анализ промышленного цикла, охватывающий все уровни данного явления и основанный на принципе восхождения от абстрактного к конкретному: от исследования товара, стоимости, прибавочной стоимости и капитала – к анализу кругооборота и оборота индивидуальных капиталов, воспроизводства совокупного общественного капитала, конкуренции, общей нормы прибыли, кредита, мирового рынка и промышленного цикла как конкретного единства всех названных отношений.

Таким образом, степень разработанности научной проблемы характеризуется, с одной стороны, наличием фундаментального теоретического наследия классиков марксизма, значительного массива советских и современных марксистских исследований, многочисленных буржуазных и эклектических концепций, а также развитого инструментария исторического, статистического и математического анализа, а с другой – отсутствием достаточно целостной теоретико-методологической системы, способной объединить изучение внутренних законов движения капитала с исследованием конкретно-исторических форм промышленного цикла. Это обусловливает необходимость критической систематизации существующих концепций, преодоления их односторонности и эклектизма, восстановления связи внешне наблюдаемых колебаний производства, инвестиций, занятости и кредита с имманентными противоречиями капиталистического воспроизводства и разработки комплексного подхода к анализу промышленного цикла на основе диалектического и историко-материалистического метода.

\subsection*{Объект исследования}
В качестве объекта настоящего исследования выступает \textit{циклическая динамика движения капиталистической экономики}.

\subsection*{Предмет исследования}
В качестве предмета настоящего исследования выступают \textit{характер, факторы, причины, закономерности и внешние проявления циклической динамики движения капиталистической экономики на современном этапе}.

\subsection*{Цель исследования}
Цель исследования заключается в теоретическом и эмпирическом обосновании промышленного цикла в качестве ключевой специфической формы диалектического развития капиталистической экономики на ее текущем конкретно-историческом этапе.

\subsection*{Задачи исследования}
Для достижения цели исследования поставлены следующие задачи:
\begin{enumerate}
    \item Определить и конкретизировать понятие цикла в экономике, проанализировать его отражение в контексте становления экономической теории; систематизировать теоретические концепции, определения, термины, методологию, подходы, инструментарий и методы его исследования;
    \item Определить структуру и характер цикла в условиях современной экономики; обозначить индикаторы, тип, уровень и исторические границы экономики, исходную сферу, фундаментальную причину, факторы, период, структуру фаз цикла, а также наличие или отсутствие явления мультицикличности в экономике; систематизировать различные теоретические концепции на основе ключевых параметров;
    \item Обосновать промышленный цикл в качестве специфической формы диалектического развития капиталистической экономики; описать его материалистическое основание, системный характер, механизм развития в контексте фундаментальных закономерностей развития капиталистической экономики как внутренне противоречивой системы;
    \item Описать процесс и динамику воспроизводства основного капитала в качестве ключевого фактора промышленного цикла в капиталистической экономике; отразить его специфический характер, инертность и периодичность воспроизводства, а также его циклообразующую роль в условиях капиталистического накопления;
    \item Осуществить конкретно-исторический анализ кризисов и общей политэкономической динамики промышленного цикла в экономике США; рассмотреть кризис перепроизводства как конкретно-историческое и социальнло-экономическое явление, описать типологию и разновидности кризисов капиталистической экономики; выделить периоды и фазы промышленного цикла в экономике США, проанализировать их продолжительность и динамику; описать особенности текущего промышленного цикла, спрогнозировать приблизительную дату и повод следующего кризиса общего перепроизводства;
    \item Провести эмпирический анализ сопутствующих индикаторов и показателей промышленного цикла в экономике США, исследовать динамику различных статистических показателей, ее изменение по фазам промышленного цикла и взаимосвязь с динамикой промышленного производства.
\end{enumerate}

\subsection*{Методы исследования}
Для целей настоящего исследования используются следующие научные методы:
\begin{itemize}
    \item \textit{диалектический и историко-материалистический методы}: составляют методологическую основу исследования, позволяя рассматривать промышленный цикл не как механическое чередование периодов роста и спада хозяйственной активности, а как исторически определенную форму движения внутренних противоречий капиталистического способа производства. Применение данного метода предполагает анализ циклической динамики в ее возникновении, развитии и изменении, раскрытие взаимосвязи производительных сил и производственных отношений, производства и обращения, накопления и потребления, капитала и наемного труда, а также выявление тех противоречий, обострение и временное разрешение которых обусловливают переход капиталистической экономики от одной фазы цикла к другой;
    \item \textit{метод восхождения от абстрактного к конкретному}: используется в качестве основного способа теоретического построения, предполагая последовательный переход от наиболее общих и простых категорий капиталистического хозяйства – товара, стоимости, денег, прибавочной стоимости и капитала – к более сложным формам его движения, включая кругооборот и оборот капитала, воспроизводство совокупного общественного капитала, конкуренцию, образование средней нормы прибыли, кредит, фиктивный капитал, мировой рынок и промышленный цикл. Такой порядок исследования позволяет представить промышленный цикл как конкретное единство многообразных экономических отношений и процессов, а не как внешнюю сумму отдельных колебаний производства, инвестиций, цен, кредита, прибыли и занятости;
    \item \textit{исторический и логический методы}: применяются в их диалектическом единстве. Исторический метод используется для исследования возникновения промышленного цикла вместе с развитием крупной машинной индустрии и мирового рынка, последовательного изменения его форм на стадиях капитализма свободной конкуренции, монополистического и государственно-монополистического капитализма, а также для периодизации кризисов и выявления конкретно-исторических особенностей отдельных циклов; логический метод позволяет воспроизвести внутреннюю структуру и необходимые связи исследуемого явления, отвлекаясь от случайных обстоятельств и располагая категории в соответствии с объективной последовательностью развития капиталистических производственных отношений;
    \item \textit{методы научной абстракции, анализа и синтеза, индукции и дедукции, обобщения и классификации}: используются при изучении научных концепций экономического цикла. Научная абстракция позволяет отделить существенные, необходимые и устойчивые связи циклического воспроизводства от случайных и второстепенных факторов; анализ применяется для расчленения промышленного цикла на его основные стороны, фазы, причины и механизмы, тогда как синтез обеспечивает их последующее соединение в целостную теоретическую конструкцию; индукция используется для обобщения фактического и статистического материала, а дедукция – для выведения конкретных форм циклической динамики из общих законов движения капитала. Метод классификации применяется при систематизации марксистских, буржуазных и эклектических концепций цикла по совокупности содержательных критериев, включая понимание источника цикличности, механизма перехода между фазами, роли кредита, инвестиций, нормы прибыли, технического прогресса и государственного регулирования;
    \item \textit{системный подход}: используется для рассмотрения промышленного цикла как целостного, внутренне расчлененного и многоуровневого процесса, охватывающего движение индивидуальных капиталов, воспроизводство совокупного общественного капитала, взаимодействие отраслей и подразделений общественного производства, национальную экономику и мировое капиталистическое хозяйство. На его основе исследуется взаимное влияние промышленной, торговой, кредитной, финансовой и внешнеэкономической составляющих цикла, а также механизм передачи кризисных импульсов между различными отраслями, секторами и странами; одновременно структурно-функциональный анализ применяется для определения места отдельных элементов капиталистического воспроизводства в общей системе и раскрытия их роли на различных фазах цикла;
    \item \textit{сравнительно-исторический метод}: применяется для сопоставления отдельных промышленных циклов, экономических кризисов и этапов развития капитализма, выявления их общих закономерностей и исторически специфических черт. Его использование позволяет разграничить устойчивое содержание циклического движения и изменяющиеся формы его проявления, установить влияние монополизации, развития кредита, государственного регулирования, интернационализации производства, финансиализации и технологических преобразований на продолжительность, глубину и структуру циклов, а также определить, какие изменения относятся к модификации механизма кризиса, а какие затрагивают более глубокие условия воспроизводства общественного капитала;
    \item \textit{статистический, динамический и структурный методы}: используются для конкретизации теоретических положений при анализе социально-экономических показателей. Исследованию подлежат динамика промышленного производства, валового накопления, инвестиций в основной капитал, загрузки производственных мощностей, занятости, заработной платы, массы и нормы прибыли, органического строения капитала, производительности труда, кредитной задолженности и других показателей, характеризующих воспроизводство и циклическое движение капитала; при этом анализ временных рядов применяется для выявления переломных моментов, последовательности фаз, продолжительности колебаний и взаимосвязи между отдельными показателями, а структурный анализ – для изучения изменений отраслевых и воспроизводственных пропорций;
    \item \textit{расчетно-аналитический и экономико-математический методы}: используются при необходимости количественной оценки политико-экономических категорий. Включает построение относительных и индексных показателей, расчет темпов роста и прироста, сглаживание динамических рядов, сопоставление средних величин, анализ пропорций и корреляционных зависимостей. Математические и эконометрические методы рассматриваются при этом не как самостоятельный источник теоретических выводов, а как вспомогательный инструмент проверки и конкретизации положений политической экономии, поскольку статистическая связь между показателями сама по себе не раскрывает причинного механизма цикла и требует содержательной интерпретации на основе анализа капиталистических производственных отношений;
    \item \textit{источниковедческий и критический методы}: применяются при изучении трудов классиков марксизма, советской и современной политико-экономической литературы, а также работ представителей немарксистских направлений. Анализ концепций осуществляется с учетом исторического контекста их формирования, исходных философских и методологических предпосылок, используемого категориального аппарата, внутренней логической непротиворечивости и соответствия фактическому материалу; особое внимание уделяется выявлению случаев, когда отдельные внешние формы циклической динамики абсолютизируются и превращаются в универсальную причину кризиса, а исследование капиталистических производственных отношений подменяется описанием процессов обращения, кредита, спроса или технических изменений.
\end{itemize}

Таким образом, применяемая совокупность методов обеспечивает соединение абстрактно-теоретического, конкретно-исторического и эмпирического уровней исследования. Диалектический и исторический материализм определяет общую методологическую направленность работы, метод восхождения от абстрактного к конкретному задает логику раскрытия предмета, исторический, сравнительный и системный методы позволяют выявить эволюцию и структуру промышленного цикла, а статистические и расчетно-аналитические инструменты служат средством конкретизации и проверки теоретических положений на материале реального движения капиталистической экономики.

\subsection*{Теоретико-методологическая база исследования}
Теоретико-методологическую базу исследования составляют положения диалектического и исторического материализма, марксистской политической экономии и теории общественного воспроизводства, позволяющие рассматривать промышленный цикл как исторически определенную форму движения капиталистического способа производства. В соответствии с данным подходом циклические колебания исследуются не как случайные отклонения от состояния экономического равновесия и не как результат воздействия отдельных внешних факторов, а как закономерное проявление внутренних противоречий производства, обращения, распределения и потребления в условиях господства капиталистических производственных отношений.

Фундаментальную теоретическую основу исследования образуют труды К. Маркса и Ф. Энгельса, в которых были раскрыты законы движения капитала, механизмы производства и реализации прибавочной стоимости, кругооборот и оборот промышленного капитала, воспроизводство совокупного общественного капитала, образование средней нормы прибыли, тенденция нормы прибыли к понижению, развитие кредита и фиктивного капитала, а также закономерный характер кризисов общего перепроизводства. Особое значение для настоящей работы имеют положения К. Маркса о противоречии между стремлением капиталистического производства к безграничному расширению и ограниченными условиями реализации произведенной стоимости, о периодическом обесценении капитала как форме восстановления нарушенных пропорций воспроизводства и о массовом обновлении основного капитала как одной из материальных основ промышленного цикла.

Существенное значение имеют труды В. И. Ленина, в которых получили дальнейшее развитие положения марксистской теории капиталистического воспроизводства, рынка, реализации общественного продукта, неравномерности экономического развития и закономерностей монополистической стадии капитализма. Ленинский анализ концентрации и централизации производства, образования монополий, слияния банковского и промышленного капитала, вывоза капитала и формирования мирового капиталистического хозяйства используется в исследовании для раскрытия исторической трансформации промышленного цикла и изменения конкретных форм кризисных процессов по мере перехода от капитализма свободной конкуренции к монополистическому капитализму. Теоретические положения И. В. Сталина об основном экономическом законе современного капитализма, общем кризисе капиталистической системы и соотношении производительных сил и производственных отношений также учитываются при анализе исторических условий изменения циклической динамики.

Важной составной частью теоретической базы выступают работы советских экономистов П. С. Максаковского, Л. А. Мендельсона, Е. С. Варги, И. А. Трахтенберга, В. Г. Фигурнова и других исследователей, развивавших марксистскую теорию кризисов и промышленного цикла. В их трудах были систематизированы представления о фазах цикла, раскрыты закономерности перехода от кризиса к депрессии, оживлению и подъему, исследованы роль обновления основного капитала, динамика инвестиций, цен, кредита, прибыли и занятости, а также разработана историческая периодизация экономических кризисов. Данные положения используются для соединения абстрактно-теоретического анализа движения капитала с конкретно-историческим исследованием циклического развития капиталистической экономики.

При анализе современных форм движения капитала учитываются труды М. Робертса, Г. Карчеди, В. Калюжного и других современных авторов, применяющих категории марксистской политической экономии к исследованию статистического материала развитых и зависимых капиталистических экономик. Особое внимание уделяется разработанным ими подходам к эмпирической оценке органического строения капитала, нормы и массы прибыли, степени эксплуатации труда, накопления основного капитала, международного перераспределения стоимости и неэквивалентного обмена. Использование данных исследований позволяет сопоставить теоретические положения о росте органического строения капитала и тенденции нормы прибыли к понижению с фактической динамикой современного капиталистического хозяйства.

Методологическая организация исследования основывается на принципе восхождения от абстрактного к конкретному. Исходным пунктом выступает анализ наиболее общих категорий капиталистического производства – товара, стоимости, денег, рабочей силы, прибавочной стоимости и капитала, после чего исследование последовательно переходит к кругообороту и обороту индивидуального капитала, воспроизводству совокупного общественного капитала, конкуренции, межотраслевому переливу капитала, образованию общей нормы прибыли, развитию кредитной системы, мирового рынка и государственно-монополистического регулирования. Промышленный цикл в результате рассматривается как конкретное единство многообразных процессов и противоречий капиталистического воспроизводства, а не как механическая совокупность колебаний отдельных макроэкономических показателей.

Важнейшим методологическим принципом работы является единство логического и исторического. Логический анализ направлен на раскрытие внутренней структуры промышленного цикла и необходимых связей между его фазами, тогда как исторический анализ позволяет исследовать изменение конкретных форм циклического движения на различных этапах развития капитализма. Такое сочетание дает возможность отделить устойчивые закономерности промышленного цикла, вытекающие из сущности капитала, от исторически преходящих форм их проявления, обусловленных монополизацией, развитием финансового капитала, государственным регулированием, интернационализацией производства, цифровизацией и изменением международного разделения труда.

Исследование опирается также на принцип противоречивости как внутреннего источника развития экономических явлений. Движение промышленного цикла анализируется через обострение и временное разрешение противоречий между производством и потреблением, производством и обращением, стоимостью и потребительной стоимостью, накоплением и реализацией, организацией производства на отдельных предприятиях и стихийностью общественного производства в целом, общественным характером труда и частнокапиталистической формой присвоения. При этом кризис понимается как насильственная форма временного восстановления единства взаимосвязанных моментов воспроизводства, относительная самостоятельность которых в ходе капиталистического накопления перерастает в их открытый разрыв.

Теоретико-методологическая база исследования включает критический анализ немарксистских концепций экономического цикла, представленных кейнсианством, монетаризмом, австрийской школой, теориями реального делового цикла, инновационными, психологическими, инвестиционными и денежно-кредитными объяснениями цикличности. Содержащиеся в данных концепциях положения о динамике инвестиций, кредита, совокупного спроса, ожиданий, технических нововведений и финансовой нестабильности рассматриваются как отражение отдельных реальных сторон капиталистического воспроизводства, однако подвергаются критической оценке с точки зрения их способности раскрыть внутреннюю необходимость промышленного цикла. Особое внимание уделяется выявлению случаев, когда отдельная форма проявления кризиса превращается в его универсальную причину, а исследование капиталистических производственных отношений подменяется анализом функциональных зависимостей между агрегированными показателями.

Таким образом, теоретико-методологическая база исследования представляет собой взаимосвязанную систему положений диалектического и исторического материализма, марксистской теории стоимости и прибавочной стоимости, теории воспроизводства общественного капитала, учения о кризисах и промышленном цикле, а также конкретно-исторических исследований развития мирового капиталистического хозяйства. Такая база обеспечивает рассмотрение промышленного цикла в единстве его сущности и форм проявления, общих закономерностей и исторической специфики, теоретического содержания и эмпирически наблюдаемой динамики, создавая основу для комплексного анализа причин, механизма, структуры и исторической эволюции циклического движения капиталистической экономики.

\subsection*{Информационно-эмпирическая база исследования}
Информационно-эмпирическую базу исследования составили официальные статистические материалы, аналитические базы данных и публикации государственных, межгосударственных и научно-исследовательских организаций, характеризующие динамику капиталистического воспроизводства, промышленного производства, накопления капитала, инвестиций, прибыли, занятости, кредита и внешнеэкономических связей. Основной массив эмпирических данных относится к экономике США, которая рассматривается как одна из наиболее развитых капиталистических экономик и занимает центральное положение в системе современного мирового хозяйства.

При определении хронологических границ отдельных промышленных циклов, дат начала и окончания рецессий, а также продолжительности фаз циклической динамики использовались материалы Национального бюро экономических исследований США – National Bureau of Economic Research, NBER. Предлагаемая NBER периодизация деловых циклов применялась как исходная эмпирическая основа для сопоставления переломных моментов экономической конъюнктуры с изменениями промышленного производства, инвестиций, нормы прибыли, занятости, загрузки производственных мощностей, кредитной активности и других показателей воспроизводства общественного капитала.

Важнейшим источником статистических данных выступили материалы Бюро экономического анализа США – Bureau of Economic Analysis, BEA. На их основе исследовались динамика валового внутреннего продукта и его компонентов, валовая и чистая добавленная стоимость, корпоративная прибыль, амортизация, валовое и чистое накопление основного капитала, инвестиции в основной капитал, доходы различных секторов экономики и иные показатели, необходимые для анализа движения стоимости, накопления капитала и воспроизводственных пропорций американской экономики.

Для формирования и обработки длительных динамических рядов использовалась база Federal Reserve Economic Data – FRED, поддерживаемая Федеральным резервным банком Сент-Луиса. Данная база применялась для получения и сопоставления показателей промышленного производства, загрузки производственных мощностей, занятости и безработицы, заработной платы, производительности труда, процентных ставок, денежного предложения, кредитной задолженности, инфляции, доходности ценных бумаг и других индикаторов, отражающих движение промышленного, торгового, ссудного и фиктивного капитала на различных фазах экономического цикла.

Дополнительными источниками эмпирической информации служили статистические материалы Федеральной резервной системы США, Бюро трудовой статистики США, Бюро переписи населения США и других американских государственных учреждений. Их данные использовались для анализа отраслевой структуры производства, динамики занятости, заработной платы и производительности труда, состояния рынка рабочей силы, изменения цен, движения промышленного капитала, инвестиционной активности и финансовых условий капиталистического воспроизводства.

При исследовании международного движения капитала, трансграничных инвестиций, деятельности транснациональных корпораций и изменения положения США в мировом хозяйстве использовались ежегодные доклады *World Investment Report*, публикуемые Конференцией ООН по торговле и развитию – UNCTAD. Содержащиеся в них данные о притоках и оттоках прямых иностранных инвестиций, накопленных объемах зарубежных активов, международных слияниях и поглощениях, глобальных цепочках создания стоимости и деятельности транснационального капитала применялись для конкретизации анализа интернационализации капиталистического производства и распространения циклических процессов в рамках мирового хозяйства.

В работе также использовались статистические и аналитические материалы международных экономических организаций, национальных статистических служб и профильных исследовательских учреждений, позволяющие проводить межстрановые и исторические сопоставления. Обращение к нескольким независимым источникам данных обеспечивало возможность проверки сопоставимости показателей, выявления различий в методологии их расчета и снижения риска односторонней интерпретации эмпирического материала.

Обработка информационно-статистической базы включала формирование временных рядов, расчет темпов роста и прироста, относительных и индексных показателей, сопоставление средних величин, выявление переломных точек и структурных изменений, а также расчет политико-экономических показателей, характеризующих накопление и движение капитала. Особое внимание уделялось сопоставимости данных, изменениям методики статистического учета, различиям между номинальными и реальными величинами, валовыми и чистыми показателями, а также разграничению текущих цен и цен постоянного базисного периода.

Официальные статистические показатели рассматривались не как непосредственное и исчерпывающее отражение сущности капиталистических отношений, а как эмпирическая форма их проявления, требующая теоретической интерпретации. Поэтому анализ данных осуществлялся на основе категориального аппарата марксистской политической экономии, что предполагало критическую оценку используемых национальными и международными организациями понятий, их сопоставление с категориями стоимости, прибавочной стоимости, постоянного и переменного капитала, органического строения капитала, массы и нормы прибыли, а также выявление ограничений буржуазной статистики при измерении действительных процессов капиталистического воспроизводства.

Таким образом, информационно-эмпирическая база исследования объединяет официальную макроэкономическую и отраслевую статистику США, материалы авторитетных научно-исследовательских баз данных, сведения о международном движении капитала и прямых иностранных инвестициях, а также результаты собственных расчетов автора. Использование данных NBER, BEA, FRED, Федеральной резервной системы, Бюро трудовой статистики США, UNCTAD и других источников позволило соединить теоретический анализ промышленного цикла с конкретным исследованием динамики капиталистической экономики и проверить основные положения работы на длительном историко-статистическом материале.

\subsection*{Научная новизна исследования}
\begin{enumerate}
    \item Произведена наиболее полная на текущий момент систематизации различных теоретических концепций цикла в экономике с применением разработанной в рамках настоящей работы авторской системы классификации на основе ключевых аспектов анализа и принципиальных дискуссионных вопросов теории цикла;
    \item Выделены и конкретизированы главные и второстепенные противоречия капиталистической экономики и ее составных элементов, лежащие в основе промышленного цикла, подробно описано их обострение и разрешение в ходе кризиса перепроизводства. Обоснована роль промышленного цикла в качестве специфической формы диалектического развития капиталистической экономики, при которой ее главные противоречия остаются неизменными, а второстепенные противоречия частично и временно снимаются насильственным путем в ходе кризиса перепроизводства, и затем воспроизводятся вновь на более высокой стадии развития капиталистической экономики;
    \item На основе положений марксистской политической экономии построена авторская модель периодического массового обновления основного капитала для случая простого воспроизводства, наглядно иллюстрирующая механизм периодизации и синхронизации моментов обновления основного капитала между различными предприятиями в капиталистической экономике;
    \item С применением конкретно-исторического и политико-экономического анализа систематизированы и классифицированы все экономические кризисы в экономике США, исходя из их датировок и типологии выделены периоды и фазы промышленного цикла в экономике США в период с 1919 года по настоящее время; рассчитана и проанализирована динамика органического строения капитала и нормы прибыли на основе актуальных статистических данных во взаимосвязи с фазами промышленного цикла, эмпирическим путем подтверждающая фундаментальные положения марксистской теории о долгосрочной тенденции повышения органического строения капитала и понижения нормы прибыли.
\end{enumerate}

\subsection*{Теоретическая значимость исследования}
Теоретическая значимость исследования заключается в развитии положений марксистской политической экономии о промышленном цикле как исторически определенной форме движения и временного разрешения внутренних противоречий капиталистического воспроизводства. В работе циклическая динамика рассматривается не как внешнее чередование периодов роста и спада и не как результат действия какого-либо единичного фактора, а как целостный процесс, возникающий из противоречивого единства производства и обращения, накопления и реализации, стоимости и потребительной стоимости, общественного характера производства и частнокапиталистической формы присвоения его результатов.

Исследование способствует теоретическому уточнению содержания категории промышленного цикла и ее разграничению с более узкими понятиями делового, инвестиционного, кредитного и финансового циклов. Промышленный цикл раскрывается как всеобщая форма движения воспроизводства общественного капитала, охватывающая динамику производства, инвестиций, занятости, прибыли, кредита, цен, обновления основного капитала и межотраслевых пропорций, тогда как отдельные разновидности циклических колебаний рассматриваются как относительно самостоятельные, но внутренне связанные моменты единого процесса капиталистического воспроизводства.

Теоретическое значение работы состоит также в систематизации существующих концепций экономического цикла на основе разграничения марксистской, буржуазной и эклектической исследовательских линий. Предлагаемый подход позволяет классифицировать теории не только по формально выделяемому источнику колебаний, но и по совокупности методологических критериев, включающих понимание природы капитала, причины кризиса, механизма смены фаз, роли нормы прибыли, кредита, технического прогресса, государственного регулирования и мирового рынка, что создает более содержательную основу для сравнительного анализа различных направлений экономической мысли.

Значимость исследования определяется применением принципа восхождения от абстрактного к конкретному к анализу циклической динамики капиталистического хозяйства. Последовательное движение от категорий товара, стоимости, прибавочной стоимости и капитала к исследованию его кругооборота и оборота, воспроизводства совокупного общественного капитала, конкуренции, средней нормы прибыли, кредита и мирового рынка позволяет раскрыть промышленный цикл как конкретное единство множества экономических процессов, а также преодолеть фрагментарность подходов, абсолютизирующих отдельные стороны кризисного движения.

Работа имеет теоретическое значение для дальнейшего исследования связи между накоплением капитала, ростом его органического строения, движением нормы прибыли, обновлением основного капитала и возникновением кризисов общего перепроизводства. Рассмотрение данных процессов в их взаимной обусловленности позволяет уточнить механизм перехода капиталистической экономики от подъема к кризису, от кризиса к депрессии и последующему оживлению, а также выявить роль обесценения капитала и восстановления условий прибыльного накопления в формировании материальной основы нового цикла.

Теоретическая значимость исследования связана также с конкретизацией положения об исторической изменчивости форм промышленного цикла. Анализ монополизации, финансиализации, развития транснационального капитала, глобальных производственных цепочек, государственного антициклического регулирования и усиления международного движения капитала позволяет показать, что современные институциональные и финансовые механизмы способны изменять продолжительность, глубину и внешние формы отдельных фаз цикла, но не устраняют его объективной основы, заключенной в противоречиях капиталистического способа производства.

Использование официальных статистических данных по экономике США и материалам международного движения капитала создает возможность конкретизации категорий марксистской политической экономии применительно к современному этапу развития капитализма. Теоретически значимым является сопоставление эмпирической динамики промышленного производства, инвестиций, корпоративной прибыли, основного капитала, занятости, кредита и прямых иностранных инвестиций с положениями о циклическом характере накопления, росте органического строения капитала и тенденции нормы прибыли к понижению, что способствует соединению абстрактно-теоретического анализа с исследованием реального движения капиталистической экономики.

Таким образом, теоретическая значимость исследования состоит в развитии комплексного политико-экономического подхода к промышленному циклу, объединяющего диалектический и историко-материалистический анализ, теорию воспроизводства общественного капитала, исследование движения нормы прибыли, кредита, основного капитала и мирового рынка. Полученные теоретические положения могут быть использованы для дальнейшей разработки марксистской теории кризисов, совершенствования классификации концепций экономического цикла, исследования современных форм капиталистической цикличности и более глубокого понимания исторических границ государственного регулирования капиталистического хозяйства.

\subsection*{Практическая значимость исследования}
Практическая значимость исследования заключается в возможности использования его результатов для анализа текущего состояния и перспектив развития капиталистической экономики, определения фазы промышленного цикла и выявления предпосылок перехода от одной его стадии к другой. Разработанный подход позволяет рассматривать изменения производства, инвестиций, прибыли, занятости, кредита и загрузки производственных мощностей не изолированно, а как взаимосвязанные элементы единого процесса воспроизводства общественного капитала.

Полученные результаты могут применяться при проведении макроэкономического и отраслевого анализа, подготовке прогнозов промышленной динамики, оценке инвестиционной активности и выявлении нарастающих воспроизводственных диспропорций. Особое практическое значение имеет возможность сопоставления движения нормы прибыли, накопления основного капитала, органического строения капитала, корпоративных доходов и кредитной задолженности с изменением темпов промышленного производства, что позволяет более обоснованно оценивать устойчивость экономического подъема и вероятность формирования кризисных тенденций.

Предложенная система показателей и методика работы с официальными статистическими данными могут использоваться при анализе экономики США и других развитых капиталистических стран. Материалы NBER, BEA, FRED, Федеральной резервной системы, Бюро трудовой статистики США и UNCTAD позволяют формировать длительные сопоставимые временные ряды, рассчитывать политико-экономические показатели и выявлять взаимосвязь между внутренним накоплением капитала, международным движением инвестиций и изменением положения национальной экономики в мировом капиталистическом хозяйстве.

Практическая значимость работы состоит также в возможности использования ее выводов для критической оценки мер государственного антициклического регулирования. Результаты исследования позволяют разграничивать меры, способные временно смягчить внешние проявления кризиса, и меры, которые лишь переносят противоречия во времени, усиливают долговую нагрузку, поддерживают обесценивающийся капитал либо перераспределяют последствия кризиса между различными классами и секторами экономики. Это создает основу для более содержательной оценки денежно-кредитной, бюджетной, инвестиционной и промышленной политики.

Разработанная классификация теорий экономического цикла может применяться в научно-исследовательской и экспертно-аналитической деятельности для систематизации различных подходов к объяснению кризисов. Она позволяет выявлять методологические основания конкретных концепций, определять границы их применимости и отличать анализ реальных механизмов циклической динамики от объяснений, сводящих кризис к отдельным внешним факторам, ошибкам регулирования или случайным колебаниям ожиданий экономических субъектов.

Материалы и выводы исследования могут использоваться в образовательном процессе при преподавании политической экономии, истории экономических учений, макроэкономики, мировой экономики и специальных дисциплин, посвященных экономическим кризисам и циклам. Разработанные теоретические положения, классификации, статистические материалы и методические подходы могут быть положены в основу учебных курсов, лекционных материалов, семинарских занятий, учебно-методических пособий и дальнейших научных исследований.

Результаты работы могут быть востребованы научными организациями, аналитическими центрами, образовательными учреждениями, профессиональными объединениями и общественными структурами, занимающимися исследованием промышленной динамики, инвестиционных процессов, международного движения капитала и социально-экономических последствий кризисов. Они позволяют оценивать циклические процессы не только с точки зрения изменения агрегированных макроэкономических показателей, но и с учетом их воздействия на занятость, заработную плату, условия труда, уровень эксплуатации и распределение созданной стоимости между трудом и капиталом.

Таким образом, практическая значимость исследования определяется возможностью применения его результатов для диагностики фаз промышленного цикла, анализа кризисных тенденций, оценки эффективности государственного регулирования, совершенствования методик обработки статистических данных и подготовки научно обоснованных прогнозов капиталистической динамики. Полученные положения могут служить инструментом комплексного анализа процессов воспроизводства общественного капитала и использоваться как в научно-исследовательской и образовательной работе, так и в экспертной оценке современных экономических процессов.

\subsection*{Область исследования}
Содержание исследования, а также основные положения и выводы соответствуют паспорту специальности 5.2.1 \textquotedbl{}Экономическая теория\textquotedbl{}:
\begin{itemize}
    \item \textit{пункт 2}: уточнение экономического содержания категории промышленного цикла, определение его структуры, внутренних противоречий и места в системе категорий воспроизводства общественного капитала, а также разграничение промышленного, делового, инвестиционного, кредитного и финансового циклов;
    \item \textit{пункт 4}: разработка и применение методологии комплексного исследования промышленного цикла, основанной на диалектическом и историко-материалистическом методе, принципе восхождения от абстрактного к конкретному и соединении теоретического, конкретно-исторического и статистического уровней анализа;
    \item \textit{пункт 5}: исследование исторического развития теорий промышленного цикла, выявление основных этапов их становления и анализ изменения представлений о причинах и механизме экономических кризисов;
    \item \textit{пункт 6}: систематизация научных школ и исследовательских программ, объясняющих природу экономических циклов, а также сравнительный анализ их предметных, методологических и теоретических оснований;
    \item \textit{пункт 7}: конкретно-историческиое исследование возникновения и эволюции промышленного цикла, периодизация кризисов капиталистической экономики и сопоставление циклической динамики на различных этапах развития капитализма;
    \item \textit{пункт 9}: исследование макроэкономических закономерностей циклического движения капиталистической экономики, включая взаимосвязь динамики производства, накопления, инвестиций, прибыли, занятости, кредита и совокупного спроса;
    \item \textit{пункт 11}: применение политико-экономического подхода к исследованию промышленного цикла, в рамках которого циклическая динамика выводится из внутренних противоречий производства, накопления и воспроизводства общественного капитала;
    \item \textit{пункт 16}: разработка теоретических подходов к исследованию промышленного цикла как формы экономических колебаний капиталистического хозяйства, раскрытие его внутренних причин, фазовой структуры, механизма движения и исторически изменяющихся форм проявления.
\end{itemize}

\subsection*{Структура исследования}
Структура исследования обусловлена поставленными целью и задачами, логикой раскрытия промышленного цикла как целостной формы движения капиталистического воспроизводства, а также применением принципа восхождения от абстрактного к конкретному. Работа включает в себя Введение, \formbytotal{totalchapter}{Глав}{}{ы}{} (разделенные на 6 параграфов), Заключение, Список использованных источников и литературы. В Приложениях приводятся дополнительные статистические материалы, расчетные таблицы, графики, методические пояснения и иные данные, конкретизирующие положения основной части работы. Работа содержит \formbytotal{TotPages}{страниц}{у}{ы}{} печатного текста, включая \formbytotal{totalcount@table}{таблиц}{у}{ы}{}, \formbytotal{totalcount@figure}{рисун}{ок}{ка}{ков} и \formbytotal{totalappendix}{приложен}{ия}{ий}{ий}.